% !TEX root = Lecture_Koszul.tex
\clearpage
\section{Minimal graded resolutions}\label{sec:minimal-resolutions}

% TODO (ChangeProposal9, Lecture 4 follow-up): After approval, refer to
% def:positive-ideal and eq:degree-zero-projection for A_+ and A/A_+ = S.
% When Ext and Tor are introduced in Lectures 6--7, explain
% Hom_A(M,S) = Hom_S(M/MA_+,S) naturally, and
% Ext_A^i(M,S) = Hom_S(Tor_i^A(M,S),S) naturally (ungraded right modules).
% For a minimal graded resolution, identify the graded Ext groups with the
% S-duals of the generator modules P_i/P_iA_+; check the shift sign carefully.
% This is not higher Hom_S(-,S), which is exact. Specify the Yoneda algebra
% action and its opposite convention when formulating Koszul duality.

\subsection{Projective covers and the graded Nakayama lemma}

Let $A=\bigoplus_{d\geq0}A_d$ be a locally finite-dimensional graded algebra with $S=A_0$ finite-dimensional semisimple.
Write
\[
A_+:=\bigoplus_{d>0}A_d.
\]
The grading makes $A_+$ a two-sided ideal, and projection onto degree $0$ identifies $A/A_+$ with $S$.
We regard $S$ as a graded right $A$-module concentrated in degree $0$, with action $s\cdot a=sa_0$ for $s\in S$ and $a=\sum_{d\geq0}a_d\in A$.

\phantomsection\label{not:chain-differentials}
All complexes below are chain complexes of graded right $A$-modules: their terms $P_i$ are indexed by homological degrees $i\geq0$, and their differentials are degree preserving maps $d_i:P_i\to P_{i-1}$ satisfying $d_{i-1}d_i=0$ for $i\geq2$.
An augmented complex to a module $M$ also has a degree preserving map $p_0:P_0\to M$ with $p_0d_1=0$.
The augmented complex is exact when $\Im d_{i+1}=\Ker d_i$ for $i\geq1$, $\Im d_1=\Ker p_0$, and $p_0$ is surjective.

\phantomsection\label{not:graded-top}
For a graded right $A$-module $M$, define its \defn{graded top} by
\[
\overline M:=M/MA_+.
\]
The quotient $\overline M$ is a graded right $S$-module.
We call $M$ \defn{bounded below} if $M_d=0$ for all sufficiently small $d$.
Every finitely generated graded module is bounded below: a finite generating set can be replaced by its finitely many homogeneous components, and $A$ has no negative degrees.

\begin{definition}\label{def:graded-projective}
A graded right $A$-module $P$ is \defn{projective} if, for every surjective degree preserving map $q:X\to Y$ and every degree preserving map $f:P\to Y$, there is a degree preserving map $\widetilde f:P\to X$ such that $q\widetilde f=f$.
\end{definition}

The shifted free module $A(-d)$ has its generator $1$ in degree $d$.
If $Q$ is a finite-dimensional projective right $S$-module, then $Q\otimes_S A(-d)$ is a projective graded right $A$-module: $Q$ is a direct summand of a finite free right $S$-module, so $Q\otimes_S A(-d)$ is a direct summand of a finite direct sum of copies of $A(-d)$.

\begin{lemma}[Graded Nakayama lemma]\label{lem:graded-nakayama}
If a bounded below graded right $A$-module $M$ satisfies $M=MA_+$, then $M=0$.
\end{lemma}

\begin{proof}
If $M\neq0$, let $d$ be the least degree with $M_d\neq0$.
Every degree $d$ element of $MA_+$ is a sum of products $m_j a_{d-j}$ with $j<d$ and $d-j>0$.
All such $m_j$ vanish by the choice of $d$, so $(MA_+)_d=0$, contradicting $M_d=(MA_+)_d$.
\end{proof}

\begin{definition}\label{def:graded-projective-cover}
A \defn{graded projective cover} of a bounded below locally finite-dimensional graded right $A$-module $M$ is a surjective degree preserving map $p:P\to M$ from a bounded below locally finite-dimensional graded projective module such that
\[
\overline p:P/PA_+\xrightarrow{\sim}M/MA_+.
\]
\end{definition}

For a surjection $p:P\to M$, the condition in Definition~\ref{def:graded-projective-cover} is equivalent to $\Ker p\subseteq PA_+$, since $p(PA_+)=MA_+$.

\begin{proposition}\label{prop:graded-projective-covers}
Every bounded below locally finite-dimensional graded right $A$-module has a graded projective cover, unique up to isomorphism over $M$.
\end{proposition}

\begin{proof}
Put $\overline M_d=(M/MA_+)_d$. Since $S$ is semisimple, the quotient map $M_d\to\overline M_d$ has a right $S$-linear section $s_d:\overline M_d\to M_d$.
Set
\[
P:=\bigoplus_{d\in\mathbb Z}(\overline M_d\otimes_S A)(-d).
\]
Each $\overline M_d$ is projective over $S$, so $P$ is a graded projective right $A$-module.
Since $M$ is bounded below and locally finite-dimensional, each degree of $P$ is a finite sum of finite-dimensional spaces; hence $P$ is locally finite-dimensional.
The maps
\[
\overline M_d\otimes_S A\longrightarrow M,\qquad
\overline m\otimes a\longmapsto s_d(\overline m)a
\]
define a degree preserving map $p:P\to M$.
On graded tops, $\overline p$ identifies each copy of $\overline M_d$ with $(M/MA_+)_d$.
Thus $\overline p$ is an isomorphism.
The quotient $C:=M/p(P)$ satisfies $C/CA_+=0$, hence $C=CA_+$; Lemma~\ref{lem:graded-nakayama} gives $C=0$.
Therefore $p$ is a projective cover.

For uniqueness, let $p:P\to M$ and $q:Q\to M$ be projective covers.
Projectivity gives a map $f:P\to Q$ with $qf=p$.
The maps $\overline p$ and $\overline q$ are isomorphisms, so $\overline f=\overline q^{-1}\overline p$ is an isomorphism.
The quotient $Q/f(P)$ has zero graded top and is bounded below, hence is zero by Lemma~\ref{lem:graded-nakayama}.
Thus $f$ is surjective; since $Q$ is projective, $f$ splits.
Its kernel is a direct summand of $P$, and its graded top is zero because $\overline f$ is injective.
Lemma~\ref{lem:graded-nakayama} gives $\Ker f=0$, so $f$ is an isomorphism over $M$.
\end{proof}

\subsection{Minimal resolutions and graded Betti numbers}

\begin{definition}\label{def:minimal-resolution}
A \defn{graded projective resolution} of $M$ is an exact sequence of degree preserving maps
\[
\cdots\longrightarrow P_2\xrightarrow{d_2}P_1\xrightarrow{d_1}P_0\xrightarrow{p_0}M\longrightarrow0
\]
with every $P_i$ graded projective.
It is \defn{minimal} if $p_0:P_0\to M$ is a graded projective cover and, for each $i\geq1$, the map $d_i$ with codomain restricted to $\Ker d_{i-1}$ is a graded projective cover, where $d_0:=p_0$.
\end{definition}

\begin{theorem}\label{thm:minimal-resolution}
Every bounded below locally finite-dimensional graded right $A$-module has a minimal graded projective resolution.
For any two such resolutions $P_\bullet\to M$ and $P'_\bullet\to M$, there are degree preserving isomorphisms $P_i\to P'_i$ commuting with the differentials and the maps to $M$.
\end{theorem}

\begin{proof}
Take a projective cover $p_0:P_0\to M$ using Proposition~\ref{prop:graded-projective-covers}, and set $K_1=\Ker p_0$.
The module $K_1$ is bounded below and locally finite-dimensional, so it has a projective cover $q_1:P_1\to K_1$.
Repeat: after defining $q_i:P_i\to K_i$, set $K_{i+1}=\Ker q_i$ and cover $K_{i+1}$ by $q_{i+1}:P_{i+1}\to K_{i+1}$.
Composing each $q_i$ with the inclusion $K_i\hookrightarrow P_{i-1}$ gives $d_i$ and an exact minimal resolution.

For uniqueness, Proposition~\ref{prop:graded-projective-covers} identifies the covers of $M$ in two minimal resolutions.
The resulting isomorphism identifies their kernels; applying the same proposition to those kernels identifies the next projective covers.
Induction gives an isomorphism in every homological degree, compatible with the differentials and the maps to $M$.
\end{proof}

Choose representatives $L_1,\ldots,L_r$ of the simple right $S$-modules, as in Definition~\ref{def:dimension-vector}.
\phantomsection\label{not:graded-betti}
For a minimal resolution $P_\bullet\to M$, define the \defn{graded Betti numbers} $\beta_{i,d,j}(M)$ by
\[
\left(P_i/P_iA_+\right)_d
\cong\bigoplus_{j=1}^r L_j^{\oplus\beta_{i,d,j}(M)},
\qquad
\beta_{i,d}(M):=(\beta_{i,d,1}(M),\ldots,\beta_{i,d,r}(M)).
\]
Here $i$ is the homological degree, $d$ is the grading degree, and $j$ records the simple right $S$-module.
Theorem~\ref{thm:minimal-resolution} shows these multiplicities do not depend on the chosen minimal resolution.
The projective cover construction also gives
\[
P_i\cong\bigoplus_{d\in\mathbb Z}\bigoplus_{j=1}^r
(L_j\otimes_S A)(-d)^{\oplus\beta_{i,d,j}(M)}.
\]
For $S=\bigoplus_{j=1}^r\Bbbk e_j$, the simple modules are $L_j=e_jS$ and $L_j\otimes_S A\cong e_jA$.

\begin{example}\label{ex:minimal-A2}
This example shows how a non-field degree zero algebra gives distinct projective covers for its different simple modules.
Let $A=\Bbbk A_2$, with $S=\Bbbk e_{11}\oplus\Bbbk e_{22}$ and $A_1=\Bbbk e_{12}$.
The simple right $S$-modules are $L_1=e_{11}S$ and $L_2=e_{22}S$, with projective lifts $Q_1=e_{11}A$ and $Q_2=e_{22}A$.
The quotient map $Q_1\to L_1$ has kernel $\Bbbk e_{12}$, while $Q_2=L_2$.
Sending the degree $1$ generator $e_{22}$ to $e_{12}$ gives the minimal resolution
\[
0\longrightarrow Q_2(-1)\xrightarrow{d_1}Q_1\longrightarrow L_1\longrightarrow0,
\qquad d_1(e_{22})=e_{12}.
\]
\end{example}

\begin{example}\label{ex:minimal-polynomial}
For $A=\Bbbk[x]$, with $\deg x=1$, the degree zero quotient module is $\Bbbk=A/(x)$.
Multiplication by $x$ gives the minimal resolution
\[
0\longrightarrow A(-1)\xrightarrow{d_1}A\xrightarrow{\pi}\Bbbk\longrightarrow0,
\qquad d_1(a)=xa,\quad \pi(f)=f(0).
\]
The map $d_1$ has degree $0$, is injective, and has image $xA=\Ker\pi\subset AA_+$.
Thus $\beta_{0,0}=1$ and $\beta_{1,1}=1$.
\end{example}

\begin{example}\label{ex:minimal-exterior}
The exterior algebra on one generator gives an infinite minimal resolution, although the algebra itself is finite-dimensional.
Let $E=\Bbbk[\epsilon]/(\epsilon^2)$ with $\deg\epsilon=1$ and $E_+=\Bbbk\epsilon$.
The degree zero quotient is $\Bbbk=E/E_+$, and
\[
\cdots\xrightarrow{d_3}E(-2)\xrightarrow{d_2}E(-1)
\xrightarrow{d_1}E\xrightarrow{\pi}\Bbbk\longrightarrow0,
\qquad d_i(a)=\epsilon a,\quad \pi(a)=a_0.
\]
Each $d_i$ is degree preserving and right $E$-linear.
Since $\Ker(\epsilon\cdot -)=\epsilon E=\Im(\epsilon\cdot -)$, the sequence is exact at every $E(-i)$ and at $E$.
Every differential has image in the preceding term times $E_+$, so the resolution is minimal and $\beta_{i,i}=1$ for all $i\geq0$.
\end{example}

\begin{advanced}
If $M$ is finitely generated, its first projective cover is finitely generated.
Without a noetherian hypothesis on $A$, later kernels need not be finitely generated, so later terms of the minimal resolution need not be finitely generated.
The bounded below locally finite-dimensional setting above still gives projective covers and a minimal resolution.
\end{advanced}

